\documentclass{article}

\usepackage{amsmath,amssymb}
\usepackage{xurl}
\usepackage[hidelinks]{hyperref}
\setlength{\emergencystretch}{2em}

% Shared commands for notes in docs/paper-gaps/.

\DeclareUnicodeCharacter{2080}{\ifmmode{}_{0}\else\textsubscript{0}\fi}
\DeclareUnicodeCharacter{2081}{\ifmmode{}_{1}\else\textsubscript{1}\fi}
\DeclareUnicodeCharacter{2082}{\ifmmode{}_{2}\else\textsubscript{2}\fi}
\DeclareUnicodeCharacter{2099}{\ifmmode{}_{n}\else\textsubscript{n}\fi}

\newcommand{\Lean}{\textsc{Lean}}
\ExplSyntaxOn
\NewDocumentCommand{\leanid}{m}{
  \ifmmode
    \textsf{\detokenize{#1}}
  \else
    \group_begin:
    \urlstyle{sf}
    \tl_set:Nn \l_tmpa_tl {#1}
    \tl_replace_all:Nnn \l_tmpa_tl {\_} {_}
    \exp_args:NV \path \l_tmpa_tl
    \group_end:
  \fi
}
\ExplSyntaxOff
\providecommand{\tr}{\operatorname{tr}}
\newcommand{\tnleanIssueBase}{https://github.com/LionSR/TNLean/issues/}
\newcommand{\tnleanPrBase}{https://github.com/LionSR/TNLean/pull/}
\newcommand{\tnleanDocBase}{https://sirui-lu.com/TNLean/docs/}
\newcommand{\ghissue}[1]{\href{\tnleanIssueBase#1}{GitHub issue \##1}}
\newcommand{\ghpr}[1]{\href{\tnleanPrBase#1}{PR \##1}}
\newcommand{\doclink}[1]{\href{\tnleanDocBase#1.html}{\path{#1}}}

% Dirac notation and the identity operator, as in blueprint/src/macros/common.tex.
% \providecommand keeps note-local definitions authoritative.
\usepackage{dsfont}
\providecommand{\ket}[1]{|#1\rangle}
\providecommand{\bra}[1]{\langle #1|}
\providecommand{\braket}[2]{\langle #1 | #2 \rangle}
\providecommand{\ketbra}[2]{|#1\rangle\!\langle #2|}
\providecommand{\Id}{\mathds{1}}
\providecommand{\one}{\mathds{1}}
\providecommand{\C}{\mathbb{C}}
\providecommand{\MN}[1]{M_{#1}(\C)}
\providecommand{\MD}{M_D(\C)}

% External referencing for paper sources.  The build helper writes sanitized
% auxiliary files under build/paper-aux/ and excludes duplicated source labels.
\usepackage{xr}

% Import a generated paper auxiliary file with a caller-supplied label prefix.
% The candidate paths support builds from docs/paper-gaps/, the repository root,
% and build/paper-aux/ itself.
\newcommand{\importpaperaux}[2]{%
  \IfFileExists{../../build/paper-aux/#2.aux}{%
    \externaldocument[#1]{../../build/paper-aux/#2}%
  }{%
    \IfFileExists{build/paper-aux/#2.aux}{%
      \externaldocument[#1]{build/paper-aux/#2}%
    }{%
      \IfFileExists{#2.aux}{%
        \externaldocument[#1]{#2}%
      }{}%
    }%
  }%
}

% General external reference macros
\newcommand{\paperref}[2]{\ref{#1-#2}}
\newcommand{\papereqref}[2]{(\ref{#1-#2})}

% CPSV16 source (1606.00608 / MPDO-22-12-17-2.tex) external document and shortcuts
\importpaperaux{cpsv16-}{1606.00608/MPDO-22-12-17-2-tnlean}
\newcommand{\cpsvref}[1]{\paperref{cpsv16}{#1}}
\newcommand{\cpsveqref}[1]{\papereqref{cpsv16}{#1}}

% Verdict marker: \gapnote{<kind>}{<status>}. Kind is the policy
% classification (clarification, local-correction, scope-restriction,
% unfaithful, false-source, open-gap); status is open, wip (elimination
% underway), resolved, or historical. Machine-read by the site index;
% prints a verdict line.
\newcommand{\gapnote}[2]{\par\noindent\textbf{Verdict:} #1 (#2).\par}


\title{Example Parent Hamiltonians in Cirac--Perez-Garcia--Schuch--Verstraete 2021}
\author{}
\date{2026-06-04}

\begin{document}
\maketitle
\gapnote{scope-restriction}{open}

\section{Purpose}

This note records the source scope of the example-level parent-Hamiltonian
targets in \href{https://github.com/LionSR/TNLean/issues/2315}{issue \#2315}.
The issue concerns the GHZ, cluster, and AKLT tensors already present in the
formal development. Its mathematical objective is to connect these tensors to
their parent Hamiltonians and corresponding ground spaces. The review does not,
however, state every standard example formula used by that objective. Such
formulas require either another citation or an explicit derivation.

\section{Statements present in the review}

For physical index $i$ and auxiliary indices $\alpha,\beta$, the review defines
the GHZ tensor by
\begin{align}
    A^i_{\alpha\beta}
        &=
        \delta_{i=\alpha=\beta}.
    \label{eq:rmp_example_parent_hamiltonian_scope_ghz_tensor}
\end{align}
Its periodic matrix product state is
\begin{align}
    \ket{\mathrm{GHZ}}
        &=
        \sum_{i=0}^{d-1}\ket{i,i,\ldots,i}.
    \label{eq:rmp_example_parent_hamiltonian_scope_ghz_state}
\end{align}
Both formulas appear in Ref.~\cite[lines~2341--2345]{Cirac2021Matrix}.

In its discussion of symmetry breaking, the review uses the ferromagnetic Ising
Hamiltonian
\begin{align}
    H_{\mathrm{Ising}}
        &=
        -\sum_i Z_iZ_{i+1}.
    \label{eq:rmp_example_parent_hamiltonian_scope_ising_hamiltonian}
\end{align}
It identifies the symmetric ground state as the superposition of the all-up and
all-down configurations and gives the two physical tensor matrices
\begin{align}
    A^\uparrow
        &=
        \ketbra{0}{0},
    \label{eq:rmp_example_parent_hamiltonian_scope_ghz_up_matrix}\\
    A^\downarrow
        &=
        \ketbra{1}{1}.
    \label{eq:rmp_example_parent_hamiltonian_scope_ghz_down_matrix}
\end{align}
These statements appear in Ref.~\cite[lines~1194--1196]{Cirac2021Matrix}. The
review later states that a GHZ parent Hamiltonian is a ferromagnetic Ising
Hamiltonian with a twofold-degenerate ground space
\cite[line~2205]{Cirac2021Matrix}. It records the two-dimensional GHZ PEPS
tensor separately \cite[lines~2417--2421]{Cirac2021Matrix}.

A local source gives the two-site GHZ support and an equivalent local
interaction:
\begin{align}
    \mathcal G_2
        &=
        \operatorname{span}
        \left\{\ket{00},\ket{11}\right\},
    \label{eq:rmp_example_parent_hamiltonian_scope_ghz_two_site_support}\\
    h_{\mathrm{GHZ}}
        &=
        \frac{1}{2}I
        -
        \left(\ketbra{00}{00}
            +
            \ketbra{11}{11}\right).
    \label{eq:rmp_example_parent_hamiltonian_scope_ghz_local_interaction}
\end{align}
Here $I$ denotes the identity on the two-site Hilbert space, and the second
expression is stated up to an additive constant.\footnote{See
\path{Papers/1210.6613/UncleHMPS.tex}:432--455.}

For the one-dimensional cluster state, the review gives the MPS matrices
\begin{align}
    A^0
        &=
        \ketbra{0}{+},
    \label{eq:rmp_example_parent_hamiltonian_scope_cluster_matrix_zero}\\
    A^1
        &=
        \ketbra{1}{-}.
    \label{eq:rmp_example_parent_hamiltonian_scope_cluster_matrix_one}
\end{align}
It also describes preparation by controlled-$Z$ gates from
$\ket{+}^{\otimes N}$ \cite[lines~2364--2369]{Cirac2021Matrix}. The
two-dimensional cluster PEPS appears separately in the examples appendix
\cite[lines~2424--2429]{Cirac2021Matrix}.

The review's one-dimensional cluster passage does not state the standard
stabilizer Hamiltonian
\begin{align}
    H_{\mathrm{cluster}}
        &=
        -\sum_i Z_{i-1}X_iZ_{i+1}.
    \label{eq:rmp_example_parent_hamiltonian_scope_cluster_stabilizer}
\end{align}
The older MPS review instead states that cluster states are unique ground states
of the three-body interactions
\begin{align}
    H_{ZXZ}
        &=
        \sum_i
        \sigma_i^z
        \sigma_{i+1}^x
        \sigma_{i+2}^z,
    \label{eq:rmp_example_parent_hamiltonian_scope_cluster_three_body}
\end{align}
and supplies an explicit two-matrix MPS representation
\cite[local TeX lines~374--387]{PerezGarcia2007MPSReps}. The formal target
should therefore use the $ZXZ$ convention for the controlled-$Z$ cluster state
unless it explicitly introduces a local-basis variant and proves the
equivalence. The general normal-MPS parent-Hamiltonian theorem gives a unique
ground state for a suitable parent Hamiltonian once normality and range
reduction have been formalized
\cite[lines~2087--2090]{Cirac2021Matrix}; it does not by itself identify the
explicit cluster stabilizer interaction.

For the one-dimensional AKLT state, the review constructs the tensor from
spin-$\tfrac12$ singlets followed by projection onto the joint spin-$1$
subspace, and gives its matrices in two bases
\cite[lines~2372--2393]{Cirac2021Matrix}. It states that the AKLT tensor has
injectivity length $L_0=2$ and that a two-site parent interaction already
suffices for a unique ground state. The stated local check is
\begin{align}
    \mathcal G_{1,2}\cap\mathcal G_{2,3}
        &=
        \mathcal G_{1,2,3}.
    \label{eq:rmp_example_parent_hamiltonian_scope_aklt_intersection}
\end{align}
This criterion appears in Ref.~\cite[line~2095]{Cirac2021Matrix}.

For open boundaries, the review also describes dangling spin-$\tfrac12$ edge
modes. It states that the parent-Hamiltonian ground-space degeneracy is at least
the square of the virtual irreducible-representation dimension and that the
AKLT case has a four-dimensional ground space containing only one spin-$0$
state \cite[lines~1169--1174]{Cirac2021Matrix}.

The local review passages do not display the standard description of the AKLT
Hamiltonian as a sum of spin-$2$ projectors. That description requires a
citation to the AKLT source or a derivation from spin addition before it can be
treated as a source-faithful target here. The local sources instead display the
polynomial Hamiltonian
\begin{align}
    H_{\mathrm{AKLT}}
        &=
        \sum_i
        \left[\vec S_i\cdot\vec S_{i+1}
            +
            \frac{1}{3}
            \left(\vec S_i\cdot\vec S_{i+1}\right)^2\right],
    \label{eq:rmp_example_parent_hamiltonian_scope_aklt_polynomial}
\end{align}
together with the unnormalised matrices
\begin{align}
    \left\{A^{-1},A^0,A^{+1}\right\}
        &=
        \left\{\sigma^z,\sqrt{2}\sigma^+,-\sqrt{2}\sigma^-\right\}.
    \label{eq:rmp_example_parent_hamiltonian_scope_aklt_matrices}
\end{align}
These formulas appear in
Ref.~\cite[local TeX lines~340--349]{PerezGarcia2007MPSReps}. The same
polynomial Hamiltonian, with an additive constant, appears in the local
Schuch--P\'erez-Garc\'ia--Cirac source.\footnote{See
\path{Papers/1010.3732/paper_v3.tex}:2126--2130.}

\section{Current formal boundary}

The formal example modules already define the tensors and prove elementary
structural properties. The GHZ tensor has matrices
\begin{align}
    A^0
        &=
        \ketbra{0}{0},
    \label{eq:rmp_example_parent_hamiltonian_scope_formal_ghz_zero}\\
    A^1
        &=
        \ketbra{1}{1}.
    \label{eq:rmp_example_parent_hamiltonian_scope_formal_ghz_one}
\end{align}
The corresponding file proves that this tensor is an RFP, is not injective, and
has the expected $\mathbb Z_2$ symmetry.\footnote{See
\path{TNLean/MPS/Examples/GHZ.lean}:54--86 and the module summary at
\path{TNLean/MPS/Examples/GHZ.lean}:21--29.}

The cluster module uses the reflected one-dimensional convention
\begin{align}
    A^0
        &=
        \ketbra{+}{0},
    \label{eq:rmp_example_parent_hamiltonian_scope_formal_cluster_zero}\\
    A^1
        &=
        \ketbra{-}{1}.
    \label{eq:rmp_example_parent_hamiltonian_scope_formal_cluster_one}
\end{align}
It proves non-injectivity at one site and injectivity after blocking two
sites.\footnote{See \path{TNLean/MPS/Examples/Cluster.lean}:62--78 and
\path{TNLean/MPS/Examples/Cluster.lean}:200--215.}

The AKLT module defines the spin-$1$ tensor and proves two-block injectivity,
hence normality.\footnote{See \path{TNLean/MPS/Examples/AKLT.lean}:49--58 and
\path{TNLean/MPS/Examples/AKLT.lean}:222--237.}

The general parent-Hamiltonian API is also available. For a block length $L$,
the local interaction is the orthogonal projector onto
$\mathcal G_L(A)^\perp$, and the finite-chain parent Hamiltonian is the sum of
its translates.\footnote{See
\path{TNLean/MPS/ParentHamiltonian/Defs.lean}:66--75 and
\path{TNLean/MPS/ParentHamiltonian/Defs.lean}:23--27.}
The formal library already proves that this parent Hamiltonian annihilates the
MPS vector.\footnote{See
\path{TNLean/MPS/ParentHamiltonian/Basic.lean}:118--128.}

The normal-MPS range-reduction theorem still depends on the periodic-boundary
comparison. This is the remaining formal gap in the general unique-ground-state
statement.\footnote{Formal declaration:
\leanid{MPSTensor.wrapped_mirror_witness_agree_of_chainGroundSpace}. See
\path{TNLean/MPS/ParentHamiltonian/UniqueGroundState.lean}:822--846 and
\path{TNLean/MPS/ParentHamiltonian/UniqueGroundState.lean}:856--894.}

\section{Source-faithful order of work}

A source-faithful formalization of issue \#2315 should proceed as follows.
\begin{enumerate}
    \item For GHZ, first compare the parent Hamiltonian with the ferromagnetic
          Ising Hamiltonian and prove the two-dimensional ground-space
          statement used by the review
          \cite[line~2205]{Cirac2021Matrix}. The two-site projector formula
          may be taken from the local uncle-Hamiltonian source.

    \item For the cluster state, first connect the existing normality proof to
          the general parent-Hamiltonian construction, or prove the three-body
          $ZXZ$ interaction stated in the older MPS review
          \cite[local TeX lines~374--387]{PerezGarcia2007MPSReps}. State any
          other convention only after proving the required local-basis or
          index-shift equivalence.

    \item For AKLT, first use the two-site criterion
          (\ref{eq:rmp_example_parent_hamiltonian_scope_aklt_intersection}) and
          the existing two-block injectivity proof. The polynomial Hamiltonian
          in
          (\ref{eq:rmp_example_parent_hamiltonian_scope_aklt_polynomial}) is
          locally sourced. Introduce the spin-$2$ projector formulation only
          after citing or proving its equivalence with the polynomial form.
          Treat the fourfold open-boundary degeneracy separately from the
          periodic unique-ground-state theorem.
\end{enumerate}

\section{Verdict}

This is a scope restriction. Issue \#2315 is source-faithful only when it
separates three types of targets: parent-Hamiltonian consequences explicitly
stated in the review, explicit stabilizer or spin-projector formulas, and
open-boundary edge-mode ground spaces. These are distinct mathematical
statements with distinct source obligations.

\bibliographystyle{alpha}
\bibliography{references}

\end{document}
